\chapter{Pattern Recognition}

\section{Vai trò của Pattern Recognition}
\defaultauthor

Pattern Recognition trong IMOVE không phải một mô hình machine learning dự đoán giao thông. Dự án nhận diện các cấu trúc và hành vi lặp lại để tái sử dụng chiến lược xử lý ổn định: validation, scheduling, ranking, context adjustment, preference learning và failure handling.

\section{Các pattern được áp dụng}
\defaultauthor

\begin{longtable}{p{2.8cm}p{4cm}p{4cm}p{4cm}}
\toprule
\textbf{Pattern} & \textbf{Mẫu nhận diện} & \textbf{Implementation} & \textbf{Ý nghĩa} \\
\midrule
\endfirsthead
\toprule
\textbf{Pattern} & \textbf{Mẫu nhận diện} & \textbf{Implementation} & \textbf{Ý nghĩa} \\
\midrule
\endhead
Validation & Place phải có key và thời gian hợp lệ & \code{_REQUIRED_KEYS}, \code{_validate_time()} & Phát hiện data lỗi lúc startup \\
Structural & Trip gồm Day; Day gồm Leg & \code{TripPlan}, \code{DayPlan}, \code{LegResponse} & Contract nhất quán \\
Scheduling & Điểm gần và phù hợp time window được xét trước & \code{_day_bucketed_greedy()} & Tạo lịch thực dụng \\
Context & Mưa và giờ cao điểm đổi preference hiệu dụng & \code{ContextSnapshot}, \code{_effective_weights()} & Recommendation thích nghi \\
Ranking & Mọi mode được so sánh trên bốn dimension & \code{score_alternatives()} & Chọn recommended mode \\
Behavioral & Thay đổi mode lặp lại thể hiện preference & \code{learn_from_implicit()} & Học xu hướng người dùng \\
Resilience & Mode/API có thể không khả dụng & \code{NoRouteError}, estimates & Graceful degradation \\
Adaptation & Alert cần được kiểm tra định kỳ & polling LTA/weather & Đề xuất thay đổi hành trình \\
\bottomrule
\end{longtable}

\section{Greedy Scheduling Pattern}
\defaultauthor

\subsection{Nhận diện bài toán}
\defaultauthor

Bài toán có nét tương đồng với nearest-neighbor routing, nhưng phải đồng thời xét số ngày, thời lượng tham quan, opening hours, best-time window và địa điểm buổi tối. Vì vậy, Planning Agent dùng heuristic \code{_day_bucketed_greedy()} thay vì exhaustive search.

\subsection{Quy trình}
\defaultauthor

\begin{figure}[H]
\centering
\begin{tikzpicture}[node distance=0.85cm]
  \node[imovenode] (a) {Classify places};
  \node[imovenode,below=of a] (b) {Create day buckets};
  \node[imovenode,below=of b] (c) {Filter feasible candidates};
  \node[imovenode,below=of c] (d) {Pick nearest candidate};
  \node[imovenode,below=of d] (e) {Update clock and position};
  \node[imovenode,below=of e] (f) {Assign evening places};
  \node[imovenode,below=of f] (g) {Return groups and warnings};
  \draw[imovearrow] (a) -- (b);
  \draw[imovearrow] (b) -- (c);
  \draw[imovearrow] (c) -- (d);
  \draw[imovearrow] (d) -- (e);
  \draw[imovearrow] (e.east) -- ++(1,0) |- (c.east);
  \draw[imovearrow] (e) -- (f);
  \draw[imovearrow] (f) -- (g);
\end{tikzpicture}
\caption{Pattern lập lịch tham lam có ràng buộc thời gian.}
\label{fig:greedy-pattern}
\end{figure}

Mỗi lần chọn cần quét tập địa điểm còn lại, nên độ phức tạp xấp xỉ \(O(n^2)\). Khoảng cách Haversine được dùng cho heuristic nội bộ \cite{haversine}; route thật được lấy sau đó từ OneMap.

\section{Context-aware Ranking Pattern}
\defaultauthor

\subsection{Relative normalization}
\defaultauthor

Mỗi alternative được mô tả bằng duration, cost, walking minutes và transfers. Vì các đại lượng khác đơn vị, hệ thống chuẩn hóa tương đối trong cùng một leg:

\[
N(x) = 1 - \frac{x - x_{\min}}{x_{\max} - x_{\min}}
\]

Nếu mọi alternative có cùng giá trị trên một dimension, dimension đó nhận giá trị trung lập \(1.0\).

\subsection{Điều chỉnh theo context}
\defaultauthor

Mưa làm giảm walking weight; giờ cao điểm tăng mức quan trọng của việc giảm chuyển tuyến. Các soft constraint như \code{minimize_walking} hoặc \code{minimize_fee} tiếp tục điều chỉnh weight, sau đó toàn bộ weight được normalize lại.

\section{Behavioral Pattern Learning}
\defaultauthor

\begin{figure}[H]
\centering
\begin{tikzpicture}[node distance=1cm and 1.4cm]
  \node[imovedata] (feedback) {Implicit feedback};
  \node[imoveagent,right=of feedback] (count) {Count supported\\mode-change patterns};
  \node[imovenode,above right=of count] (mrt) {\(\geq 2\) BUS to MRT};
  \node[imovenode,below right=of count] (walk) {\(\geq 2\) changes to WALK};
  \node[imovedata,right=2cm of count] (prefs) {User preferences};
  \draw[imovearrow] (feedback) -- (count);
  \draw[imovearrow] (count) -- (mrt);
  \draw[imovearrow] (count) -- (walk);
  \draw[imovearrow] (mrt) -- (prefs);
  \draw[imovearrow] (walk) -- (prefs);
\end{tikzpicture}
\caption{Memory Agent chuyển hành vi lặp lại thành preference có cấu trúc.}
\label{fig:memory-pattern}
\end{figure}

\section{Resilience Pattern}
\defaultauthor

OneMap alternatives được lấy độc lập; một mode thất bại không làm hỏng toàn bộ request. Hệ thống có thể tạo Grab estimate khi drive route không khả dụng, tạo walking estimate cho quãng ngắn và chỉ raise \code{NoRouteError} khi không còn phương án sử dụng được. Đây là graceful degradation, không phải dự đoán lỡ chuyến.
